% ════════════════════════════════════════════
% 具体分析（5.1.1）
% ════════════════════════════════════════════
\subsubsection{诊断思路}

本部分属于数据理解与分布诊断，采用描述统计、相关分析与方差分解相结合的方法。这一部分不追求拟合模型，而是要把“数据能支持什么样的模型”说清楚，为后续变量取舍提供依据。

首先对原始文件的 8 个字段逐列识别数据类型、取值范围、唯一值数与缺失情况，判断各列应作为连续变量、名义变量还是有序变量处理。接着根据记录的组织方式判定数据粒度与结构，明确一行代表一条帖子、样本来自同一主页且只有一次观测，从而确认其为帖子级截面数据。然后按发帖前能否被运营人员决定，把变量划分为发帖决策变量、账号状态变量与结果变量，并说明各自的解释边界。最后计算曝光量的偏度、峰度与分位数，计算粉丝规模与月份的 Pearson 相关、Spearman 相关、月份对粉丝规模对数的解释度以及月内方差占比，用可量化的证据判断二者能否同时进入回归。具体分析流程如图\ref{fig:analysis1_flow}所示。

\begin{figure}[ht]
    \begin{center}
        \begin{tikzpicture}[x=1cm, y=0.7cm]
            \node[box] (n11) at (0, 0) {读取原始数据};
            \node[box] (n12) at (5, 0) {识别变量类型};
            \node[box] (n13) at (10, 0) {判定数据粒度};
            \node[box] (n22) at (5, -3) {划分因果角色};
            \node[box] (n21) at (0, -3) {判定数据结构};
            \node[box] (n23) at (10, -3) {统计分布特征};
            \node[box] (n31) at (0, -6) {计算相关结构};
            \node[box] (n32) at (5, -6) {量化共线证据};
            \node[box] (n33) at (10, -6) {输出理解结论};
            \draw[arrow] (n11) -- (n12);
            \draw[arrow] (n12) -- (n13);
            \draw[arrow] (n23) -- (n22);
            \draw[arrow] (n22) -- (n21);
            \draw[arrow] (n13) -- ++(0,-1.0) -| (n23);
            \draw[arrow] (n21) -- (n31);
            \draw[arrow] (n31) -- (n32);
            \draw[arrow] (n32) -- (n33);
        \end{tikzpicture}
        \caption{数据诊断流程图}
        \label{fig:analysis1_flow}
    \end{center}
\end{figure}
